\section{Materials and Methods}
\label{sec:methods}

\subsection{Dataset and provenance}
\label{sec:dataset}

The local dataset contained 1,030 rows, eight input variables, and one target variable, compressive strength in MPa. The input variables were cement, blast-furnace slag, fly ash, water, superplasticizer, coarse aggregate, fine aggregate, and age. This variable set follows the intended concrete-strength benchmark context described by the UCI repository \citep{Yeh1998_UCIConcreteDataset}; however, the generation log for the current project labels the local source as \texttt{DETERMINISTIC\_FALLBACK\_NOT\_UCI}. Consequently, all descriptive statistics and model outputs below are treated as fallback demonstration outputs.

\begingroup
\catcode`\_=12
\begin{table}
\caption{Dataset summary statistics. Mixture variables are kg/m3, age is days, and compressive strength is MPa. Source label: DETERMINISTIC_FALLBACK_NOT_UCI.}
\label{tab:dataset-summary}
\begin{tabular}{lrrrrrr}
\toprule
Variable & N & Mean & SD & Min & Median & Max \\
\midrule
cement & 1030 & 318.244 & 127.229 & 102.213 & 317.053 & 539.367 \\
blast\_furnace\_slag & 1030 & 102.084 & 118.639 & 0.000 & 39.397 & 358.927 \\
fly\_ash & 1030 & 53.574 & 65.159 & 0.000 & 0.000 & 199.959 \\
water & 1030 & 186.071 & 35.418 & 122.312 & 187.755 & 246.885 \\
superplasticizer & 1030 & 11.085 & 10.742 & 0.000 & 9.471 & 31.988 \\
coarse\_aggregate & 1030 & 975.166 & 97.390 & 801.274 & 975.549 & 1144.902 \\
fine\_aggregate & 1030 & 788.998 & 115.824 & 594.201 & 787.164 & 992.940 \\
age & 1030 & 52.710 & 78.089 & 1.000 & 28.000 & 365.000 \\
compressive\_strength & 1030 & 29.518 & 18.190 & 2.300 & 30.907 & 73.941 \\
\bottomrule
\end{tabular}
\end{table}

\endgroup

The generated summary table reported no missing numeric values in the local data. The compressive-strength target had a mean of 29.518~MPa, standard deviation of 18.190~MPa, minimum of 2.300~MPa, and maximum of 73.941~MPa. These values are not used as evidence about the official UCI dataset unless independent provenance verification is completed.

\subsection{Train/test split and preprocessing}
\label{sec:split_preprocessing}

The artifact-generation pipeline used an 80/20 split, yielding 824 training rows and 206 test rows. Numeric predictors were standardized for the ordinary least-squares and ridge-regression models. The training-mean baseline predicted the training-set mean response for all test observations. No claim is made that the split is unique or optimal; it is a fixed demonstration split used to produce the accompanying tables and figures.

\subsection{Ridge model}
\label{sec:ridge_model}

Ridge regression was used as the selected model because it provides a transparent regularized baseline. For standardized predictors, the fitted coefficients minimize
\begin{equation}
  \sum_{i=1}^{n}\left(y_i - \beta_0 - \mathbf{x}_i^\top\boldsymbol{\beta}\right)^2 + \alpha \sum_{j=1}^{p}\beta_j^2,
\end{equation}
where \(y_i\) is compressive strength, \(\mathbf{x}_i\) is the standardized predictor vector, \(p=8\), and \(\alpha\) is the ridge penalty. The selected model used \(\alpha = 1\). Ordinary least squares with \(\alpha=0\) was retained as a closely related comparator, and the training-mean baseline was included to check that the fitted models improved on a trivial prediction rule.

\subsection{Cross-validation and evaluation metrics}
\label{sec:evaluation}

The ridge penalty was selected from a small candidate grid using five-fold cross-validation on the training data, a standard validation approach for predictive modelling \citep{Stone1974_CrossValidation}. Performance was reported as root mean square error (RMSE), mean absolute error (MAE), coefficient of determination (\(R^2\)), and prediction bias in MPa. Test-set metrics were computed once on the held-out test set.

\subsection{Diagnostic analyses}
\label{sec:diagnostics}

The generated diagnostic outputs included a dataset overview, parity plot, residual plot, standardized coefficient plot, and age-response display. A single-feature ablation analysis refitted the selected ridge workflow after removing one input variable at a time. The ablation table is interpreted as a model-dependence diagnostic for the local fallback data, not as a causal analysis of concrete chemistry.
